\documentclass[11pt]{article}

\usepackage[margin=1in]{geometry}
\usepackage[T1]{fontenc}
\usepackage[utf8]{inputenc}
\usepackage{lmodern}
\usepackage{hyperref}
\usepackage{enumitem}
\usepackage{booktabs}

\title{VoxelMiner: Modular Real-Time Graphics Engine and Procedural Generation Experiments}
\author{VoxelMiner Repository Write-up}
\date{\today}

\begin{document}
\maketitle

\begin{abstract}
VoxelMiner is a modular graphics codebase centered on a WebGL2 rendering engine and a set of demonstration scenes, with an additional C++ procedural generation sample using Magnum/Corrade. The project emphasizes reusable scene architecture, interactive camera systems, configurable lighting, and practical experimentation through focused examples. This write-up summarizes repository structure, implementation details, design trade-offs, and future work opportunities.
\end{abstract}

\section{Project Overview}

The repository combines two complementary tracks:
\begin{enumerate}[leftmargin=*]
  \item A JavaScript/WebGL2 engine in \texttt{lib/} with reusable rendering abstractions.
  \item Example applications in \texttt{example/}, including both web demos and a native C++ procedural-generation prototype.
\end{enumerate}

The guiding objective is to provide an educational-yet-extensible graphics foundation where scene composition, camera behavior, and shading can be modified quickly without tightly coupling feature logic.

\section{Repository Structure}

\begin{itemize}[leftmargin=*]
  \item \texttt{lib/}: Core engine modules (camera, renderer, scene graph, shader model, primitives, texture and utility helpers).
  \item \texttt{example/basic/}: Baseline runtime integration with GUI controls and input management.
  \item \texttt{example/scan/}: Demonstrates rendering of a high-vertex-count statue asset.
  \item \texttt{example/car/}: Vehicle-focused example variant.
  \item \texttt{example/cloth-sim/}: Cloth-simulation themed prototype.
  \item \texttt{example/cpp-gen/}: C++/Magnum procedural house generator and interactive viewer.
  \item \texttt{README.md}: Public-facing feature summary.
  \item \texttt{bib.md}: Attribution and external resource references.
\end{itemize}

\section{JavaScript Engine Architecture}

\subsection{Core Modules}

\begin{description}[leftmargin=0cm,labelindent=0cm]
  \item[\texttt{lib/index.js}] Public entry point exporting reusable engine components.
  \item[\texttt{lib/scene.js}] \texttt{SceneGraph} stores shape list, point lights, lighting coefficients, and optional overlays (skybox/crosshair).
  \item[\texttt{lib/camera.js}] Camera model supporting FPS and arcball interaction modes with movement and projection updates.
  \item[\texttt{lib/renderer.js}] \texttt{WebGLRenderer} orchestrates render order: skybox, lit scene, then crosshair overlay.
  \item[\texttt{lib/shader-models/phong-shader.js}] Main lighting path using Phong shading and texture mode switches.
\end{description}

\subsection{Rendering Flow}

Each frame follows a clear sequence:
\begin{enumerate}[leftmargin=*]
  \item Clear color/depth buffers.
  \item Render skybox if present.
  \item Bind and configure Phong shader uniforms from scene and camera state.
  \item Render scene geometry list.
  \item Render crosshair overlay if enabled.
\end{enumerate}

This pass ordering simplifies background handling and overlay composition while keeping the main geometry path unified.

\subsection{Lighting and Shading}

The shader model includes:
\begin{itemize}[leftmargin=*]
  \item Ambient, diffuse, and specular terms with adjustable scalar factors.
  \item Up to four point lights via a fixed-size shader array.
  \item Shape-conditional behavior for specialized falloff-based rendering.
  \item Texture-source routing and diagnostic visualization modes.
\end{itemize}

\section{Input and Interaction Model}

The basic demo wires controls through an input manager and browser pointer lock:
\begin{itemize}[leftmargin=*]
  \item Keyboard movement in world space.
  \item Mouse-driven look/orbit panning.
  \item Wheel input for forward/backward movement or arcball radius changes.
  \item Runtime parameter controls through \texttt{lil-gui} for camera and lighting tuning.
\end{itemize}

This design allows immediate visual feedback when adjusting rendering parameters and camera behavior.

\section{C++ Procedural Generation Example}

\subsection{Objective}

The C++ sample in \texttt{example/cpp-gen/} demonstrates procedural house generation using Magnum's rendering stack.

\subsection{Implementation Notes}

\begin{itemize}[leftmargin=*]
  \item Procedural primitives are assembled through helper functions for triangles, quads, boxes, and roof sections.
  \item Randomized parameters control width, depth, wall height, roof height, and door dimensions.
  \item A Phong shader pipeline consumes position/normal/color vertex data.
  \item Real-time controls support movement (WASD) and regeneration (Space).
\end{itemize}

\subsection{Build Flow}

The provided script:
\begin{enumerate}[leftmargin=*]
  \item Configures CMake with a Debug build.
  \item Uses \texttt{CMAKE\_PREFIX\_PATH} (defaulting to \texttt{/opt/homebrew}) for dependency discovery.
  \item Builds executable \texttt{MyApplication}.
\end{enumerate}

\section{Design Trade-offs}

\begin{itemize}[leftmargin=*]
  \item \textbf{Simplicity vs. scalability:} fixed maximum light count and straightforward render pass ordering keep complexity low, at the cost of advanced scalability features.
  \item \textbf{Flexibility vs. strict abstraction:} examples can quickly prototype behavior, though some feature implementations remain demo-centric.
  \item \textbf{Feature breadth vs. completion depth:} several capabilities are represented at prototype level (e.g., collision hooks) rather than full production systems.
\end{itemize}

\section{Current Limitations}

\begin{itemize}[leftmargin=*]
  \item Camera collision checks are currently stubbed.
  \item Some README-listed roadmap features are partial or not yet fully integrated.
  \item Performance-oriented paths such as instancing are not fully realized.
\end{itemize}

\section{Future Work}

\begin{itemize}[leftmargin=*]
  \item Implement robust collision and world interaction mechanics.
  \item Add automated regression checks for camera behavior and shader parameter correctness.
  \item Expand shader variants and support richer material systems.
  \item Strengthen example documentation with per-scene goals and control summaries.
\end{itemize}

\section{Attribution Notes}

The repository includes external assets and references (e.g., skyboxes, textures, model data, educational materials), documented in \texttt{bib.md}. Those attributions should be preserved in derivative documentation and distribution artifacts.

\end{document}
